% ---------------------------------------------------------------------------
% Measuring a Pauli seen as +X by the graph state at v=2, on the path 1-2-3-4,
% pivoting on the min-degree neighbour b=1 (node a in the macros).
% Top row    : standard graph-only rule -- tau_b tau_v tau_b, then delete v.
% Bottom row : framed graph state -- R_b, R_v, then delete v (outcome +1,
%              no canonicalisation), with the vertex basis printed at every
%              vertex and the dual basis of v (in [...]) stacked on top:
%              T_v is its second letter.  Accent-coloured
%              bases are the ones that have left (+Z,+X): these are the side
%              effects the frame carries and a bare graph cannot.
% gold = edge added at this step, red dashed = edge removed at this step.
% Each row is scaled on its own so the two panel sizes stay legible; the
% headings and byproduct lines sit outside the boxes and keep their font.
% ---------------------------------------------------------------------------
{\small
% basis labels: \bl unchanged, \ba carrying a side effect (same style, no highlight)
\newcommand{\bl}[1]{{\large\color{Vcol!75}$#1$}}
\newcommand{\ba}[1]{\bl{#1}}
% bare panel (no frame data): #1 = edges
\newcommand{\sqg}[1]{%
  \begin{tikzpicture}[baseline={([yshift=-3.6pt]0,0.5)},x=1.6cm,y=1.35cm,gv/.append style={font=\large,minimum size=8mm},
      dele/.style={Xcol,line width=2pt,dashed}]
    \node[xv] (v) at (0,1) {$2$};
    \node[gv] (a) at (0,0) {$1$};
    \node[gv] (b) at (1,1) {$3$};
    \node[gv] (c) at (1,0) {$4$};
    #1
  \end{tikzpicture}}
\newcommand{\sqgd}[1]{%
  \begin{tikzpicture}[baseline={([yshift=-3.6pt]0,0.5)},x=1.6cm,y=1.35cm,gv/.append style={font=\large,minimum size=8mm}]
    \node[gv] (a) at (0,0) {$1$};
    \node[gv] (b) at (1,1) {$3$};
    \node[gv] (c) at (1,0) {$4$};
    #1
  \end{tikzpicture}}
% framed panel: #1..#4 = bases at v,a,b,c ; #5 = edges ; #6 = tag beside v ;
%               #7 = dual basis of v (stacked above its vertex basis)
\newcommand{\sqgf}[7]{%
  \begin{tikzpicture}[baseline={([yshift=-3.6pt]0,0.5)},x=1.8cm,y=1.35cm,gv/.append style={font=\large,minimum size=8mm},
      dele/.style={Xcol,line width=2pt,dashed}]
    \node[xv,label={[label distance=-2pt,align=center]above:#7\\[3pt]#1}] (v) at (0,1) {$2$};
    \node[gv,label={[label distance=-2pt]below:#2}] (a) at (0,0) {$1$};
    \node[gv,label={[label distance=-2pt]above:#3}] (b) at (1,1) {$3$};
    \node[gv,label={[label distance=-2pt]below:#4}] (c) at (1,0) {$4$};
    #5
    \node[font=\large,text=Xcol,anchor=south] at (0,2.2) {#6};
  \end{tikzpicture}}
\newcommand{\sqgfd}[4]{%
  \begin{tikzpicture}[baseline={([yshift=-3.6pt]0,0.5)},x=1.8cm,y=1.35cm,gv/.append style={font=\large,minimum size=8mm}]
    \node[gv,label={[label distance=-2pt]below:#1}] (a) at (0,0) {$1$};
    \node[gv,label={[label distance=-2pt]above:#2}] (b) at (1,1) {$3$};
    \node[gv,label={[label distance=-2pt]below:#3}] (c) at (1,0) {$4$};
    #4
  \end{tikzpicture}}
% arrow with enlarged label
\newcommand{\bigarrowlbl}[2][]{{\Large$\xrightarrow[\ \textstyle #1\ ]{\ \textstyle #2\ }$}}
% the graphs, shared by both rows
\newcommand{\ezero}{\draw[ge](a)--(v);\draw[ge](v)--(b);\draw[ge](b)--(c);}
% tau_1 on the leaf 1 leaves the edges unchanged
\newcommand{\eone} {\ezero}
\newcommand{\etwo} {\draw[ge](a)--(v);\draw[ge](v)--(b);\draw[ge](b)--(c);
                    \draw[ge](a)--(b);}
\newcommand{\ethree}{\draw[dele](v)--(b);\draw[ge](a)--(v);\draw[ge](b)--(c);
                     \draw[ge](a)--(b);}
\newcommand{\efinal}{\draw[ge](a)--(b);\draw[ge](b)--(c);}
%
\setlength{\tabcolsep}{2pt}
% \noindent{\color{Vcol}\textbf{Graph state}}\par\medskip

% graph-state row (a) removed: never referenced on its own
% \noindent\parbox[c]{\dimexpr\linewidth-2em}{\resizebox{\linewidth}{!}{%
% \begin{tabular}{@{}c@{}c@{}c@{}c@{}c@{}c@{}c@{}c@{}c@{}}
% \sqg{\ezero} & \bigarrowlbl{\localComp{1}} & \sqg{\eone}
%   & \bigarrowlbl{\localComp{2}} & \sqg{\etwo}
%   & \bigarrowlbl{\localComp{1}} & \sqg{\ethree}
%   & \bigarrowlbl{\mathrm{del}\,2} & \sqgd{\efinal}
% \end{tabular}}}\hfill\parbox[c]{1.6em}{\centering(a)}\par\smallskip

% % \noindent{\color{Vcol}\footnotesize byproduct: untracked}\par
% \bigskip

% \noindent{\color{Vcol}\textbf{Eulerian frame} }\par\medskip

\noindent\resizebox{\linewidth}{!}{%
\begin{tabular}{@{}c@{}c@{}c@{}c@{}c@{}c@{}c@{}}
\sqgf{\bl{(Z,X)}}{\bl{(Z,X)}}{\bl{(Z,X)}}{\bl{(Z,X)}}{\ezero}{$\transPauli_2=+X$}%
      {\bl{[Z,X]}}
  & \bigarrowlbl{\reframing{1}}
  & \sqgf{\ba{(Z,Y)}}{\ba{(Y,X)}}{\bl{(Z,X)}}{\bl{(Z,X)}}{\eone}{$\transPauli_2=-Y$}%
      {\ba{[Z,-Y]}}
  & \bigarrowlbl{\reframing{2}}
  & \sqgf{\ba{(-X,Y)}}{\ba{(Y,-Z)}}{\ba{(Z,Y)}}{\bl{(Z,X)}}{\etwo}{$\transPauli_2=-Z$}%
      {\ba{[-Y,-Z]}}
  & \bigarrowlbl{\mathrm{del}\,2}
  & \sqgfd{\ba{(Y,Z)}}{\ba{(Z,-Y)}}{\bl{(Z,X)}}{\efinal}
\end{tabular}}\par\smallskip

% \noindent{\color{Vcol}\footnotesize byproduct: $\zeta=1$ flips $\wC{}$ at $a$, $b$ and $c$}%
}
